\DocumentMetadata{
  pdfversion=2.0,
  pdfstandard=ua-2,
  lang=en-US,
  tagging=on,
  tagging-setup={math/setup=mathml-SE,tabsorder=structure}
}
\documentclass[11pt,letterpaper]{article}
\usepackage[margin=0.68in,includefoot]{geometry}
\usepackage{newcomputermodern}
\usepackage{amsmath,amscd,mathtools}
\providecommand{\square}{\mdlgwhtsquare}
\usepackage[hidelinks]{hyperref}
\hypersetup{
  pdftitle={Algebra PhD Exam, January 2014},
  pdfauthor={University of Florida Department of Mathematics},
  pdfsubject={Graduate mathematics examination},
  pdfdisplaydoctitle=true,
  bookmarksopen=true
}
\setcounter{secnumdepth}{0}
\setlength{\parindent}{0pt}
\setlength{\parskip}{0.28em}
\setlength{\emergencystretch}{3em}
\setlength{\leftmargini}{2.15em}
\setlength{\leftmarginii}{2.65em}
\setlength{\leftmarginiii}{3em}
\renewcommand{\labelenumi}{\textbf{\theenumi.}}
\pagestyle{empty}
\raggedbottom
\allowdisplaybreaks
\newenvironment{examproblems}[1]{%
  \begin{enumerate}%
  \setcounter{enumi}{#1}%
  \setlength{\topsep}{0.35em}%
  \setlength{\partopsep}{0pt}%
  \setlength{\itemsep}{0.48em}%
  \setlength{\parsep}{0.18em}%
}{\end{enumerate}}
\newenvironment{examsubparts}{%
  \begin{enumerate}%
  \setlength{\topsep}{0.18em}%
  \setlength{\partopsep}{0pt}%
  \setlength{\itemsep}{0.16em}%
  \setlength{\parsep}{0.1em}%
}{\end{enumerate}}
\newenvironment{examinstructions}{%
  \par\begingroup\itshape
}{%
  \par\endgroup\vspace{0.18em}
}
\makeatletter
\renewcommand\section{\@startsection{section}{1}{0pt}%
  {-1.25ex plus -.25ex minus -.1ex}%
  {0.55ex plus .1ex}%
  {\normalfont\large\bfseries}}
\renewcommand{\@maketitle}{%
  \newpage\null\vskip -1.2em%
  {\footnotesize\bfseries
    \href{https://www.ufl.edu/}{University of Florida}, \href{https://math.ufl.edu/}{Department of Mathematics}\par}%
  \vskip 0.18em%
  \tagstructbegin{tag=Title}%
    {\LARGE\bfseries\href{https://gma.math.ufl.edu/exams/algebra-phd/}{Algebra PhD Exam}, January 2014\par}%
  \tagstructend%
  \vskip 0.35em%
  \tagmcbegin{artifact}%
    \hrule height 1pt%
  \tagmcend%
  \vskip 0.55em%
}
\makeatother
\title{Algebra PhD Exam, January 2014}
\author{}
\date{}
\begin{document}
\maketitle
\thispagestyle{empty}
\begin{examinstructions}
Answer seven problems. (If you turn in more, the first seven will be graded.) Within reason, you may quote theorems as long as you state them clearly.
\end{examinstructions}
\begin{examinstructions}
Below, ring means associative ring with identity, and module means unital module.
\end{examinstructions}
\begin{examproblems}{0}
\item
Let \(f(x)=\sum_{i=0}^{n} a_i x^i\in\mathbb{Z}[x]\) be a polynomial with integer coefficients of degree \(n>1\). Suppose that for some~\(k\), with \(0<k<n\), and some prime~\(p\), we have \(p\nmid a_n\); \(p\nmid a_k\); \(p\mid a_i\) for \(i=0,\ldots,k-1\); and \(p^2\nmid a_0\). Show that \(f(x)\) has a factor of degree at least~\(k\) which is irreducible in \(\mathbb{Z}[x]\).
\item
Calculate the Galois group of \(x^5-12x+2\) over~\(\mathbb{Q}\), the field of rational numbers. Justify your answer carefully.
\item
This problem has two parts.
\begin{examsubparts}
\item[\textnormal{(a)}]
Give an example of fields \(M\supset L\supset K\) such that \(M/L\) and \(L/K\) are normal extensions, but \(M/K\) is not normal. Justify your answer.
\item[\textnormal{(b)}]
Let \(L\supset K\) be fields such that \([L:K]=8\). Prove that there exist \(\alpha,\beta,\gamma\in L\) such that \(L=K(\alpha,\beta,\gamma)\).
\end{examsubparts}
\item
Prove that the group defined by generators~\(a\) and~\(b\) and relations \(a^2=b^3=1\), \(abab=1\) is isomorphic to \(S_3\).
\item
Let~\(R\) be a ring with 1 and let 
\[
0\to A\xrightarrow{f}B\xrightarrow{g}C\to 0
\]
 be a short exact sequence of (left)~\(R\)-modules. Show that for any (right)~\(R\)-module~\(M\), the sequence 
\[
M\otimes_R A\xrightarrow{1_M\otimes f}M\otimes_R B\xrightarrow{1_M\otimes g}M\otimes_R C\to 0
\]
 of tensor products is exact.
\item
State and prove Hilbert’s Basis Theorem.
\item
Prove the Lying-Over Theorem: Let~\(S\) be an integral extension of an integral domain~\(R\), and let~\(P\) be a prime ideal of~\(R\). Then, there exists a prime ideal~\(Q\) of~\(S\), such that \(Q\cap R=P\).
\item
Prove that a commutative ring~\(R\) with identity is local if and only if, for all \(r,s\in R\), \(r+s=1_R\) implies that either~\(r\) or~\(s\) is a unit.
\item
Let~\(R\) be a ring.
\begin{examsubparts}
\item[\textnormal{(a)}]
Define what it means for an~\(R\)-module~\(M\) to be projective.
\item[\textnormal{(b)}]
Prove that any free~\(R\)-module is projective.
\item[\textnormal{(c)}]
Prove that there exists some ring~\(R\) and some projective~\(R\)-module~\(P\) such that~\(P\) is a projective module but not a free module.
\end{examsubparts}
\item
Determine up to isomorphism all semisimple noncommutative rings with \(512=2^9\) elements.
\item
Let~\(C\) be the category of all finitely generated~\(\mathbb{Z}\)-modules, where~\(\mathbb{Z}\) is the ring of integers.
\begin{examsubparts}
\item[\textnormal{(a)}]
Recall the definition of (direct) product in an arbitrary category~\(D\).
\item[\textnormal{(b)}]
Prove from your definition that, given a finite set~\(S\) of objects in~\(C\), there is a product~\(P\) of~\(S\) in~\(C\).
\item[\textnormal{(c)}]
Prove from your definition that there exists some set~\(T\) of objects in~\(C\), such that there is no product of~\(T\) in~\(C\).
\end{examsubparts}
\end{examproblems}
\end{document}
